% ==============================================================================
% CHAPITRE 9 : PROBABILITÉS (VERSION ÉLÈVE)
% ==============================================================================
\chapter{Probabilités \& Expériences Aléatoires}

% ------------------------------------------------------------------------------
% 1. ACTIVITÉ D'INVESTIGATION (SITUATION PROFESSIONNELLE : COMMERCE / LOTERIE)
% ------------------------------------------------------------------------------
\begin{activite}{La Roue de la Chance Commerciale (Vente \& Animation Magasin)}
Pour l'inauguration d'un nouveau magasin d'outillage, les clients peuvent faire tourner une roue de loterie partagée en \textbf{8 secteurs d'aires égales} :
\begin{itemize}[leftmargin=1.5em, itemsep=2pt]
    \item \textbf{4 secteurs \og Goodies \fg{}} (porte-clés, mètre ruban ou casquette de chantier) ;
    \item \textbf{2 secteurs \og Bon d'achat de 10 \euro{} \fg{}} ;
    \item \textbf{1 secteur \og Bon d'achat de 50 \euro{} \fg{}} ;
    \item \textbf{1 secteur \og Perdu \fg{}}.
\end{itemize}

\vspace{2mm}
\noindent
\begin{minipage}[c]{0.44\linewidth}
\centering
\begin{tikzpicture}[scale=0.95]
    % Disque 8 secteurs
    \draw[fill=sky!20, draw=slate800, line width=0.8pt] (0,0) -- (0:2.1cm) arc (0:45:2.1cm) -- cycle;
    \node[font=\bfseries\scriptsize, slate900] at (22.5:1.4cm) {Goodies};
    
    \draw[fill=amber!25, draw=slate800, line width=0.8pt] (0,0) -- (45:2.1cm) arc (45:90:2.1cm) -- cycle;
    \node[font=\bfseries\scriptsize, amber!90!black] at (67.5:1.4cm) {10 \euro{}};
    
    \draw[fill=sky!20, draw=slate800, line width=0.8pt] (0,0) -- (90:2.1cm) arc (90:135:2.1cm) -- cycle;
    \node[font=\bfseries\scriptsize, slate900] at (112.5:1.4cm) {Goodies};
    
    \draw[fill=emerald!25, draw=slate800, line width=0.8pt] (0,0) -- (135:2.1cm) arc (135:180:2.1cm) -- cycle;
    \node[font=\bfseries\scriptsize, emerald!90!black] at (157.5:1.4cm) {50 \euro{}};
    
    \draw[fill=sky!20, draw=slate800, line width=0.8pt] (0,0) -- (180:2.1cm) arc (180:225:2.1cm) -- cycle;
    \node[font=\bfseries\scriptsize, slate900] at (202.5:1.4cm) {Goodies};
    
    \draw[fill=rose!25, draw=slate800, line width=0.8pt] (0,0) -- (225:2.1cm) arc (225:270:2.1cm) -- cycle;
    \node[font=\bfseries\scriptsize, rose!90!black] at (247.5:1.4cm) {Perdu};
    
    \draw[fill=sky!20, draw=slate800, line width=0.8pt] (0,0) -- (270:2.1cm) arc (270:315:2.1cm) -- cycle;
    \node[font=\bfseries\scriptsize, slate900] at (292.5:1.4cm) {Goodies};
    
    \draw[fill=amber!25, draw=slate800, line width=0.8pt] (0,0) -- (315:2.1cm) arc (315:360:2.1cm) -- cycle;
    \node[font=\bfseries\scriptsize, amber!90!black] at (337.5:1.4cm) {10 \euro{}};

    % Contour et centre
    \draw[slate900, line width=1.5pt] (0,0) circle (2.1cm);
    \fill[slate900] (0,0) circle (4pt);
    
    % Flèche repère
    \draw[fill=colSecondary, draw=slate900, line width=1pt] (0,2.45) -- (-0.25,2.15) -- (0.25,2.15) -- cycle;
\end{tikzpicture}
\end{minipage}\hfill
\begin{minipage}[c]{0.54\linewidth}
\begin{tcolorbox}[colback=slate100, colframe=colPrimary, arc=5pt, boxrule=0.8pt, top=6pt, bottom=6pt, left=8pt, right=8pt]
\sffamily\small
\textbf{Problématique :}
\begin{enumerate}[leftmargin=1.2em, itemsep=4pt]
    \item Combien d'issues possibles compte cette expérience aléatoire ?
    \item Quelle fraction simplifiée représente la chance de gagner un bon d'achat de $10\text{ \euro{}}$ ?
    \item Quelle est la probabilité (en pourcentage) de repartir avec un lot (bon d'achat ou goodies) ?
\end{enumerate}
\end{tcolorbox}
\end{minipage}

\vspace{3mm}
\noindent\textbf{Espace de recherche :}
\lignesreponse[5]
\end{activite}

\newpage

% ------------------------------------------------------------------------------
% 2. COURS : VOCABULAIRE DES PROBABILITÉS & ÉVÉNEMENTS
% ------------------------------------------------------------------------------
\section{Vocabulaire des Probabilités \& Événements}

\begin{cadredef}[Notions Fondamentales de Probabilité]
\begin{itemize}[leftmargin=1.5em, itemsep=2pt]
    \item Une \textbf{expérience aléatoire} est une expérience dont le résultat dépend du hasard. Chaque résultat possible est une \textbf{issue}.
    \item La \textbf{probabilité} d'un événement $A$, notée $P(A)$, est un nombre compris entre \textbf{0} et \textbf{1} (ou entre $0\text{ \%}$ et $100\text{ \%}$) :
    \[ \mathbf{0 \le P(A) \le 1} \]
    \item \textbf{Événement certain} : $P = \mathbf{1}$ \quad $\bullet$ \quad \textbf{Événement impossible} : $P = \mathbf{0}$.
    \item \textbf{Événement contraire $\bar{A}$} (l'événement inverse) : $\mathbf{P(\bar{A}) = 1 - P(A)}$.
\end{itemize}
\end{cadredef}

\begin{cadreexemple}
Si la probabilité qu'une commande présente un retard est $P(R) = 0{,}08$ ($8\text{ \%}$), alors la probabilité que la commande arrive à l'heure (événement contraire $\bar{R}$) est :
\[ P(\bar{R}) = 1 - P(R) = 1 - \pointille[1.5cm] = \pointille[1.5cm] \qquad (\text{soit } \pointille[1.5cm]\text{ \%}) \]
\end{cadreexemple}

% ------------------------------------------------------------------------------
% 3. COURS : SITUATION D'ÉQUIPROBABILITÉ & ARBRES PONDÉRÉS
% ------------------------------------------------------------------------------
\section{Loi d'Équiprobabilité \& Arbre des Possibles}

\begin{cadredef}[Loi d'Équiprobabilité \& Règle de l'Arbre]
\begin{itemize}[leftmargin=1.5em, itemsep=2pt]
    \item \textbf{Équiprobabilité} (issues ayant toutes la même chance) :
    \[ \mathbf{P(A) = \frac{\text{Nombre d'issues favorables à } A}{\text{Nombre total d'issues possibles}}} \]
    \item \textbf{Arbre pondéré à 2 épreuves} : la somme des branches issues d'un n\oe ud vaut \textbf{1}. Pour trouver la probabilité au bout d'un chemin, on \textbf{multiplie} les probabilités.
\end{itemize}
\end{cadredef}

\vspace{1mm}
\begin{center}
\begin{tikzpicture}[scale=0.9, >=latex]
    % Racine
    \coordinate (O) at (0, 0);
    \fill[slate800] (O) circle (3pt);

    % Niveau 1 : Face et Pile
    \coordinate (F) at (2.8, 1.2);
    \coordinate (P) at (2.8, -1.2);

    \draw[slate700, line width=0.9pt] (O) -- (F) node[midway, above left=1pt, font=\small] {$\frac{1}{2}$};
    \draw[slate700, line width=0.9pt] (O) -- (P) node[midway, below left=1pt, font=\small] {$\frac{1}{2}$};

    \node[draw=colPrimary, fill=colPrimary!10, rounded corners=3pt, inner sep=4pt, font=\sffamily\bfseries\scriptsize] at (F) {Face ($F$)};
    \node[draw=colPrimary, fill=colPrimary!10, rounded corners=3pt, inner sep=4pt, font=\sffamily\bfseries\scriptsize] at (P) {Pile ($P$)};

    % Niveau 2 depuis Face
    \coordinate (FG) at (6.8, 1.8);
    \coordinate (FR) at (6.8, 0.6);

    \draw[slate700, line width=0.9pt] (F) -- (FG) node[midway, above=1pt, font=\small] {$\frac{1}{6}$};
    \draw[slate700, line width=0.9pt] (F) -- (FR) node[midway, below=1pt, font=\small] {$\frac{5}{6}$};

    \node[draw=emerald!80!black, fill=emerald!15, rounded corners=3pt, inner sep=3pt, font=\sffamily\bfseries\scriptsize] at (FG) {Gagné ($G$)};
    \node[draw=slate500, fill=slate100, rounded corners=3pt, inner sep=3pt, font=\sffamily\bfseries\scriptsize] at (FR) {Perdu ($R$)};

    % Niveau 2 depuis Pile
    \coordinate (PG) at (6.8, -0.6);
    \coordinate (PR) at (6.8, -1.8);

    \draw[slate700, line width=0.9pt] (P) -- (PG) node[midway, above=1pt, font=\small] {$\frac{1}{6}$};
    \draw[slate700, line width=0.9pt] (P) -- (PR) node[midway, below=1pt, font=\small] {$\frac{5}{6}$};

    \node[draw=emerald!80!black, fill=emerald!15, rounded corners=3pt, inner sep=3pt, font=\sffamily\bfseries\scriptsize] at (PG) {Gagné ($G$)};
    \node[draw=slate500, fill=slate100, rounded corners=3pt, inner sep=3pt, font=\sffamily\bfseries\scriptsize] at (PR) {Perdu ($R$)};
\end{tikzpicture}
\end{center}

\begin{cadreretenu}[L'Essentiel du Chapitre 9]
\begin{itemize}[leftmargin=1.5em, itemsep=2pt]
    \item Une probabilité est toujours comprise entre $\mathbf{0}$ et $\mathbf{1}$ (jamais négative, jamais supérieure à $1$).
    \item Équiprobabilité : $P(A) = \frac{\text{issues favorables}}{\text{issues totales}}$.
    \item Événement contraire : $P(\bar{A}) = 1 - P(A)$.
    \item Arbre pondéré : on multiplie les probabilités le long de chaque chemin.
\end{itemize}
\end{cadreretenu}

\newpage

% ------------------------------------------------------------------------------
% 4. QUESTIONS FLASH / AUTOMATISMES
% ------------------------------------------------------------------------------
\section{Questions Flash / Automatismes}

\begin{automatisme}[8 Questions Flash d'Entraînement]
\begin{enumerate}[label=\textbf{Q\arabic*.}]
    \item On lance un dé à 6 faces non truqué. Probabilité d'obtenir un nombre pair : $P = \pointille[3cm]$
    \item Dans un sac contenant 3 boules rouges et 7 boules vertes, calculer $P(\text{rouge}) = \pointille[3cm]$
    \item Exprimer en pourcentage la probabilité de tirer une boule verte de ce sac : \pointille[3cm]
    \item Si $P(E) = 0{,}35$, calculer la probabilité de l'événement contraire : $P(\bar{E}) = \pointille[3cm]$
    \item Une expérience comporte 4 issues équiprobables. Quelle est la probabilité de chaque issue ? \pointille[3cm]
    \item On lance une pièce équilibrée. Quelle est la probabilité d'obtenir Face ? \pointille[3cm]
    \item Quelle est la valeur de la probabilité d'un événement certain ? \pointille[3cm]
    \item Une probabilité peut-elle être égale à $1{,}25$ ? Justifier : \pointille[3cm]
\end{enumerate}
\end{automatisme}

% ------------------------------------------------------------------------------
% 5. TRAVAUX DIRIGÉS (TD) - EXERCICES D'ENTRAÎNEMENT
% ------------------------------------------------------------------------------
\section{Travaux Dirigés (TD) : Exercices d'Entraînement}

\begin{exercice}[2 pts]{\capRealiser}{Tirage dans un présentoir (Pâtisserie / Traiteur)}
Un artisan pâtissier prépare une boîte de \textbf{24 macarons} comprenant :
$6$ macarons chocolat, $8$ macarons vanille, $4$ macarons café et $6$ macarons framboise.
Un client choisit un macaron au hasard.
\begin{enumerate}[leftmargin=1.5em]
    \item Calculer la probabilité que le macaron soit à la vanille (fraction irréductible).
    \item Calculer la probabilité que le macaron soit au chocolat ou au café.
    \item Calculer la probabilité de l'événement : \og Le macaron n'est pas à la framboise \fg{}.
\end{enumerate}
\lignesreponse[3]
\end{exercice}

\begin{exercice}[3 pts]{\capSApproprier}{Tirage dans un jeu de 32 cartes (Jeux / Loisirs)}
On tire au hasard une carte dans un jeu de 32 cartes standard (c\oe ur, carreau, pique, trèfle) :
\begin{enumerate}[leftmargin=1.5em]
    \item Calculer la probabilité de tirer un As : $P(\text{As}) = \pointille[3cm]$
    \item Calculer la probabilité de tirer un C\oe ur : $P(\text{C\oe ur}) = \pointille[3cm]$
    \item Calculer la probabilité de tirer une carte rouge (C\oe ur ou Carreau) : $P(\text{Rouge}) = \pointille[3cm]$
    \item Calculer la probabilité de tirer une figure (Valet, Dame, Roi) : $P(\text{Figure}) = \pointille[3cm]$
\end{enumerate}
\lignesreponse[2]
\end{exercice}

\newpage

\begin{exercice}[3 pts]{\capRealiser}{Expérience à deux épreuves (Lancer combiné)}
Un jeu d'arcade consiste à lancer une pièce équilibrée (Pile ou Face), puis à tirer un jeton dans un sac contenant \textbf{1 jeton Vert} et \textbf{3 jetons Noirs}.
\begin{enumerate}[leftmargin=1.5em]
    \item Compléter l'arbre des possibles ci-dessous en inscrivant les probabilités sur chaque branche.
    \item Calculer la probabilité de l'issue \og Pile et Jeton Vert \fg{} : 
    \[ P(P \text{ et } V) = \pointille[1.8cm] \times \pointille[1.8cm] = \pointille[1.8cm] \]
\end{enumerate}

\vspace{2mm}
\begin{center}
\begin{tikzpicture}[scale=0.9, >=latex]
    % Racine
    \coordinate (O) at (0, 0);
    \fill[slate800] (O) circle (3pt);

    % Niveau 1 : Pile et Face
    \coordinate (P) at (2.8, 1.2);
    \coordinate (F) at (2.8, -1.2);

    \draw[slate700, line width=0.9pt] (O) -- (P) node[midway, above left=1pt, font=\small] {\pointille[1cm]};
    \draw[slate700, line width=0.9pt] (O) -- (F) node[midway, below left=1pt, font=\small] {\pointille[1cm]};

    \node[draw=colPrimary, fill=colPrimary!10, rounded corners=3pt, inner sep=4pt, font=\sffamily\bfseries\scriptsize] at (P) {Pile ($P$)};
    \node[draw=colPrimary, fill=colPrimary!10, rounded corners=3pt, inner sep=4pt, font=\sffamily\bfseries\scriptsize] at (F) {Face ($F$)};

    % Niveau 2 depuis Pile
    \coordinate (PV) at (6.8, 1.8);
    \coordinate (PN) at (6.8, 0.6);

    \draw[slate700, line width=0.9pt] (P) -- (PV) node[midway, above=1pt, font=\small] {\pointille[1cm]};
    \draw[slate700, line width=0.9pt] (P) -- (PN) node[midway, below=1pt, font=\small] {\pointille[1cm]};

    \node[draw=emerald!80!black, fill=emerald!15, rounded corners=3pt, inner sep=3pt, font=\sffamily\bfseries\scriptsize] at (PV) {Vert ($V$)};
    \node[draw=slate600, fill=slate100, rounded corners=3pt, inner sep=3pt, font=\sffamily\bfseries\scriptsize] at (PN) {Noir ($N$)};

    % Niveau 2 depuis Face
    \coordinate (FV) at (6.8, -0.6);
    \coordinate (FN) at (6.8, -1.8);

    \draw[slate700, line width=0.9pt] (F) -- (FV) node[midway, above=1pt, font=\small] {\pointille[1cm]};
    \draw[slate700, line width=0.9pt] (F) -- (FN) node[midway, below=1pt, font=\small] {\pointille[1cm]};

    \node[draw=emerald!80!black, fill=emerald!15, rounded corners=3pt, inner sep=3pt, font=\sffamily\bfseries\scriptsize] at (FV) {Vert ($V$)};
    \node[draw=slate600, fill=slate100, rounded corners=3pt, inner sep=3pt, font=\sffamily\bfseries\scriptsize] at (FN) {Noir ($N$)};
\end{tikzpicture}
\end{center}
\end{exercice}

\begin{exercice}[4 pts]{\capAnalyser}{Contrôle de Qualité de Composants (Électronique / Industrie)}
Une unité de production fabrique des modules électroniques pour le secteur automobile. Sur un échantillon de \textbf{500 modules} testés en fin de chaîne :
\begin{itemize}[leftmargin=1.5em, itemsep=1pt]
    \item $465$ modules ne présentent aucun défaut (entièrement conformes) ;
    \item $25$ modules présentent un défaut mineur de soudure réparable ;
    \item $10$ modules présentent un composant grillé (non réparables / rebutés).
\end{itemize}

\begin{enumerate}[leftmargin=1.5em]
    \item On prélève un module au hasard. Calculer la probabilité qu'il soit conforme : $P(C) = \pointille[3cm]$
    \item Calculer la probabilité que le module soit réparable : $P(R) = \pointille[3cm]$
    \item Calculer le pourcentage de modules rebutés par l'usine.
    \item L'usine s'est fixé un objectif qualité d'au moins $92\text{ \%}$ de modules conformes du premier coup. Cet objectif est-il atteint ? Justifier avec un calcul précis.
\end{enumerate}
\lignesreponse[4]
\end{exercice}

% ------------------------------------------------------------------------------
% 6. TP TICE / CALCULATRICE NUMWORKS
% ------------------------------------------------------------------------------
\section{TP TICE : Simuler le Hasard sur NumWorks}

\begin{tpinfo}[Calculatrice NumWorks]{Générer des Nombres Aléatoires \& Tirages}
\begin{enumerate}[leftmargin=1.5em, itemsep=2pt]
    \item \textbf{Lancer un dé à 6 faces (Application Calculs)} :
    \begin{itemize}
        \item Ouvrir l'application \textbf{Calculs}.
        \item Appuyer sur la touche \textbf{Boîte à outils} (\faTools\ Toolbox) $\rightarrow$ sélectionner \textbf{Probabilités} $\rightarrow$ \textbf{Aléatoire} $\rightarrow$ \texttt{randint(1,6)}.
        \item Appuyer plusieurs fois sur la touche \textbf{EXE} pour observer des tirages aléatoires successifs.
    \end{itemize}
    \item \textbf{Simuler une loi de probabilité (Application Probabilités)} :
    \begin{itemize}
        \item Ouvrir l'application \textbf{Probabilités} sur l'écran d'accueil.
        \item Sélectionner la loi \textbf{Uniforme} (discrète) pour visualiser la répartition théorique équiprobable.
    \end{itemize}
\end{enumerate}
\end{tpinfo}

\newpage

% ------------------------------------------------------------------------------
% 7. VERS LE BREVET (DNB PRO)
% ------------------------------------------------------------------------------
\section{Vers le Brevet (DNB Pro)}

\begin{sujetdnb}[DNB Pro Métropole][10 pts]{La Tombola du Club Omnisports}
Pour financer l'achat de matériel d'entraînement, un club omnisports organise une grande tombola. L'urne contient \textbf{100 billets} scellés, identiques et indiscernables au toucher :
\begin{itemize}[leftmargin=1.5em, itemsep=1pt]
    \item 1 billet \og Gros Lot \fg{} (un VTT d'une valeur de 350 \euro{}) ;
    \item 4 billets \og Lot Sport \fg{} (un sac de sport complet avec ballon) ;
    \item 15 billets \og Lot Consolation \fg{} (une gourde isotherme inox) ;
    \item Tous les autres billets sont \og Perdants \fg{}.
\end{itemize}

\begin{enumerate}[leftmargin=1.5em, itemsep=3pt]
    \item \capSApproprier\ Déterminer le nombre exact de billets perdants présents dans l'urne. \pointille[3cm]
    \item \capRealiser\ Un participant achète un billet au hasard. Calculer la probabilité des événements suivants (donner les résultats sous forme de fraction simplifiée puis en pourcentage) :
    \begin{enumerate}[label=\alph*), leftmargin=1.5em, itemsep=2pt]
        \item Événement $A$ : \og Gagner le Gros Lot \fg{} $\implies P(A) = \pointille[3cm]$
        \item Événement $B$ : \og Gagner un Lot Sport \fg{} $\implies P(B) = \pointille[3cm]$
        \item Événement $G$ : \og Gagner un lot (quel qu'il soit) \fg{} $\implies P(G) = \pointille[3cm]$
    \end{enumerate}
    \item \capAnalyser\ Définir par une phrase en français l'événement contraire $\bar{G}$ et calculer sa probabilité $P(\bar{G})$.
    \item \capValider\ Un joueur hésite à acheter un billet et affirme : \og J'ai plus de $75\text{ \%}$ de chances de ne rien gagner du tout \fg{}. A-t-il raison ? Justifier votre réponse par un calcul clair et argumenté.
\end{enumerate}

\vspace{2mm}
\noindent\textbf{Rédigez votre démarche ci-dessous :}
\lignesreponse[8]
\end{sujetdnb}
